\section{Methods}

\paragraph{Model.} The agent is a convolutional JEPA trained end-to-end from raw pixels. A squeeze-and-excitation residual convolutional front-end (no variational autoencoder) produces per-location features; a V-JEPA prototype-softmax self-distillation auxiliary objective provides a representational learning signal; a cross-attention memory reader (the cue-gated, top-down ``crossattn1'' module) attends from query tokens over spatial memory keys; a spatial xLSTM with memory width $d_{\mathrm{mem}}=128$ maintains recurrent state; and a distributional actor-critic head produces a per-frame \texttt{wait}/\texttt{change} decision. For the bottom-up control we use an element-wise-affine variant of the same architecture in which the memory reader is self-attention rather than cue-gated cross-attention.

\paragraph{Task.} The agent performs spatially-cued, value-directed orientation-change detection. A coloured cue marks one location and signals a reward magnitude by colour (cue colour-values $[1,3,5]$); a validity proportion sets the cue's spatial reliability. Each episode has seven logical frames; a change occurs on half of trials, always at frame~5, and is therefore detectable from frame~5 (f5) onward, with f6 one frame later. A trial is \emph{valid} if the change occurs at the cued location and \emph{invalid} if it occurs at an uncued location.

\paragraph{Set sizes.} SS4 (task vda4): four stimuli on a $2\times2$ grid, $25\times25$-pixel patches in a $50\times50$ image. SS9 (task vda9): nine stimuli on a $3\times3$ grid in a $75\times75$ image, with the cue at the top-left (P0) or bottom-right (P8) location. The architecture, auxiliary objective, and decision head are identical across set sizes; only the set size changes.

\paragraph{Change-lock metric.} For each frame we read the cross-attention distribution over the spatial memory keys (averaged over queries) and compute the attention mass on the key corresponding to the changed location. We report \emph{excess change-lock}: this mass minus the uniform baseline, which is $1/4=0.25$ at SS4 and $1/9=0.111$ at SS9. Reporting excess-over-uniform renders the two set sizes comparable. Values are reported at f5 (the change onset) and f6 (one frame later), separately for valid and invalid trials and for each cue placement (S1/S4 at SS4; P0/P8 at SS9). For the affine control the same metric is computed over the self-attention keys with the SS9 uniform baseline of $0.111$.

\paragraph{Psychometric protocol.} The SS4 valid-versus-invalid psychometric is measured by sweeping the orientation-change magnitude over $[0,4,8,16,28,44,64]$ degrees and recording $P(\text{detect})$ separately for valid and invalid trials. SS4 behavioural summary statistics (correct, hit, correct-rejection, false-alarm rates, reaction time, and the breakdown by cue colour-value) are computed over 600 greedy episodes.

\paragraph{Snapshots and scope.} Analyses use model snapshots at the following training iterations: SS4 at iteration 9599; SS9 cross-attention at iteration ${\approx}4600$; SS9 affine control at a comparable iteration (${\approx}4400$). The SS9 cross-attention model is a snapshot of a run still under continued training. Accordingly, the SS9 attentional results are reported as a mid-training snapshot, and the SS9 \emph{behavioural} psychometric is \textbf{not yet measured}: the large behavioural invalid-cue cost implied by the abolished SS9 change-lock is stated throughout as a prediction to be confirmed at convergence, never as a measured result. All change-lock quantities are reported as excess over the corresponding uniform baseline.
